\section{问题四：混合干扰源的半圆覆盖与自适应清除}
\subsection{定向反馈的联合约束}
设定向源方向为单位向量 $n$。在检测点 $S$ 收到信号的条件为
\begin{equation}
\norm{S-G}\le\rho,\qquad n^\mathsf T(S-G)\ge0.
\end{equation}
定向无信号的含义是
\begin{equation}
\norm{S-G}>\rho\quad\lor\quad n^\mathsf T(S-G)<0.\label{eq:q4negative}
\end{equation}
所以第三问的1000米无线电排除圆不能直接用于混合场景。可以在 $(G,\rho,n)$ 的联合状态中保留该析取，再投影至位置空间；最终实现以正观测的楔形和1500米圆盘维护保守位置域，避免把定向背面误判成距离过远。光学清除与朝向无关，因此20米失败清除圆仍可直接使用。

\subsection{任意朝向发现的局部凸包判据}
令有限检测点集为 $P$，对任意位置 $x\in\Omega$ 定义近邻集合 $P_x=P\cap\disk(x,1000)$。保证未知方向、未知半径的定向源被发现，等价于
\begin{equation}
\boxed{\forall x\in\Omega,\quad x\in\operatorname{conv}(P_x).}\label{eq:localhull}
\end{equation}
\begin{proof}
若 $x$ 是近邻点的凸组合，而某方向 $n$ 对所有近邻均满足 $n^\mathsf T(q-x)<0$，将这些不等式作同一凸组合便得到 $0<0$，矛盾。因此每个闭半圆必有一个可接收检测点。反之，若 $x$ 不属于非空有限近邻集的凸包，严格分离定理给出使全部近邻处于背面的方向；若近邻集为空则显然不能发现。
\end{proof}
该判据同时适用于全向源，且不需要事先识别源的类型。与问题三相比，仅有圆盘并集覆盖不足以满足未知朝向的要求。

\subsection{24点联合优化布局及连续单元证书}
最终 \texttt{q4-v18} 继承24个离线优化检测点，坐标见支撑材料 \texttt{strategy\_v17.py}。优化目标将开放路长与停点检测成本结合：
\begin{equation}
\min_{P,\pi}\ \frac{L(P,\pi)}5+\sum_i(5m_i+s_i),
\quad\text{满足式}\eqref{eq:localhull},
\end{equation}
其中 $m_i$、$s_i$ 是离线阶段对检测数和切频数的估计，实际在线成本由反馈决定。归档布局生成采用软最小值罚函数、L-BFGS-B坐标搜索、几何修复、可行性微调与访问顺序搜索，交替调整坐标和顺序。它是非凸启发式优化，不宣称全局最优，也不将其称为已实现的精确二阶锥规划。

24点的给定离线开放路线为18507.5米。在线仍与清除任务联合调度，故实际路线通常不同，这个数字只用于展示发现布局。

用目标圆包围方形的递归四分构造连续单元证书。对与目标圆相交的方格 $C$，取
\begin{equation}
P_C=\{q\in P:\max_{v\in\operatorname{Vert}(C)}\norm{q-v}\le1000-\tau\}.
\end{equation}
若四个顶点均在 $\operatorname{conv}(P_C)$ 内，则由凸性，整个格内任意 $x$ 都在该凸包中；同时 $P_C\subseteq P_x$，因此整个单元满足式\eqref{eq:localhull}。未通过的单元继续四分；细分到最小边长仍未通过则报告未决，不能当作已覆盖。

现有证书使用 $\tau=10^{-7}$ 米、最小边长0.02米，对24点布局检查3107个相交单元，1992个直接通过，未决单元为0。检查数含中间父单元，不能将3107解释为互不相交叶格数。这是带向内余量的浮点连续区域数值证书，未采用区间算术形式化验证；有限朝向抽样仅作交叉检验。

\subsection{在线精定位与清除}
扫描层保留所有待访问布局点，对未发现频道继续扫描。目标被发现后加入清除任务，滚动最近邻与2-opt统一安排扫描和清除。当前第四问调度继承其独立父类，并未复用第三问的or-opt开关。

精定位先移动至估计点尝试清除，失败后检查最多3点的覆盖计划。若区域仍宽，沿最近一次有效示向的射线分段推进以缩短距离；默认区域直径超过100米时启用，最多12步，单步上限450米。新的有效读数更新保守区域，近距反馈立即清除。无信号可能来自背面，射线上的距离区间仅用于局部搜索调度，不得作为真源的可靠上界。

补测、射线搜索及半径15米六点环形试清均属于启发式。中心加15米六点环只能保证覆盖至约31.53米的估计误差，无法覆盖任意35至40米误差。因此最终在这些流程失败后，针对原保守可行域构造25米格覆盖兜底，并使用所有仍保留的20米失败圆排除整格。由问题三的覆盖命题，该兜底与源朝向无关。

\subsection{结束条件与异常解释}
混合场景完成要求：所有已发现源均收到清除成功反馈，且每个未发现频道实际检测过的点集满足局部凸包覆盖条件。若完整访问24点布局且未知频道均被扫描，后一个条件由离线证书承担；成功清除的频道无需继续检测。

当前实现仍受64格、预算及重试次数上限限制。若计划因上限无法生成，或完整执行后仍未命中，分别表示覆盖流程未完成或保守域前提需要检查，不能把无信号简单解释为目标已消失。第四问的主要改进是将“发现覆盖”和“清除覆盖”分开论证，使波束背面造成的精定位困难最终由与方向无关的光学覆盖处理。
